%! Polynomial Functions and Complex Zeros — Unit 2, Lesson 2.4 — Slides (Beamer)
\documentclass[aspectratio=169, 11pt]{beamer}
\usepackage{precalculus-beamer}   % provides \navyheader{} and \sectionlabel[color]{}

\begin{document}

% TITLE SLIDE
{
\setbeamercolor{background canvas}{bg=navy}
\begin{frame}[plain]
  \vspace{2.8cm}\hspace{0.55cm}%
  \begin{minipage}{0.9\textwidth}
    {\color{goldacc}\bfseries\small LESSON 1.5}\\[0.5cm]
    {\color{white}\bfseries\LARGE Polynomial Functions and Complex Zeros}\\[1.2cm]
    {\color{skymid}\small Precalculus~$\cdot$~Unit 2: Polynomial Functions}
  \end{minipage}
\end{frame}
}

% HOOK
\begin{frame}[t, plain]
  \navyheader{Hook: Cross or Bounce?}
  \vspace{0.3cm}
  Consider the polynomial $p(x) = (x - 2)^2(x + 1)$.
  \vspace{0.3cm}
  \begin{block}{Think About It}
    \begin{itemize}
      \item Where does the graph \textbf{cross} the $x$-axis?
      \item Where does it \textbf{touch and bounce} without crossing?
      \item What is it about the \textbf{exponent} on each factor that tells you?
    \end{itemize}
  \end{block}
  \vspace{0.3cm}
  The exponent on a factor is called the \textbf{multiplicity} of that zero.
\end{frame}

% ZEROS, LINEAR FACTORS, MULTIPLICITY
\begin{frame}[t, plain]
  \navyheader{Zeros, Linear Factors, and Multiplicity}
  \vspace{0.2cm}
  \begin{block}{Factor–Zero Connection}
    If $p(a) = 0$, then $a$ is a \textbf{zero} of $p$ and $(x - a)$ is a linear factor.
  \end{block}
  \vspace{0.2cm}
  \begin{block}{Multiplicity}
    If $(x - a)$ is repeated $n$ times in the factored form, the zero has \textbf{multiplicity $n$}.\\[6pt]
    A degree-$n$ polynomial has exactly $\mathbf{n}$ complex zeros (counting multiplicities).
  \end{block}
  \vspace{0.2cm}
  \textbf{Example:} $p(x) = (x-2)^2(x+1)$
  \begin{itemize}
    \item $x = 2$ is a zero of \textbf{multiplicity 2}
    \item $x = -1$ is a zero of \textbf{multiplicity 1}
    \item Degree = $2 + 1 = 3$ \quad $\Rightarrow$ \quad exactly 3 zeros total
  \end{itemize}
\end{frame}

% GRAPH BEHAVIOR AT ZEROS
\begin{frame}[t, plain]
  \navyheader{Graph Behavior at Real Zeros}
  \vspace{0.3cm}
  \renewcommand{\arraystretch}{1.7}
  \begin{tabular}{|l|l|l|}
  \hline
  \textbf{Multiplicity} & \textbf{Graph behavior} & \textbf{Key word} \\
  \hline
  Odd (1, 3, 5, \ldots) & Graph \textbf{crosses} the $x$-axis & ``crosses'' \\
  \hline
  Even (2, 4, 6, \ldots) & Graph is \textbf{tangent} (touches, bounces) & ``bounce'' \\
  \hline
  \end{tabular}
  \vspace{0.4cm}
  \begin{block}{Example: $p(x) = (x+2)(x-1)^2(x-3)$}
    \begin{itemize}
      \item $x = -2$: multiplicity 1 (odd) $\Rightarrow$ \textbf{crosses}
      \item $x = 1$: multiplicity 2 (even) $\Rightarrow$ \textbf{tangent / bounces}
      \item $x = 3$: multiplicity 1 (odd) $\Rightarrow$ \textbf{crosses}
    \end{itemize}
  \end{block}
\end{frame}

% COMPLEX ZEROS AND CONJUGATE PAIRS
\begin{frame}[t, plain]
  \navyheader{Complex Zeros and Conjugate Pairs}
  \vspace{0.2cm}
  \begin{block}{Complex Conjugate Pairs Theorem}
    If $p(x)$ has \textbf{real coefficients} and $a + bi$ (with $b \neq 0$) is a zero,
    then $a - bi$ is also a zero.
  \end{block}
  \vspace{0.25cm}
  \textbf{Example:} $p(x)$ has real coefficients; $2 + 3i$ is a zero.
  \begin{itemize}
    \item Conjugate partner: $2 - 3i$ must also be a zero.
    \item Quadratic factor from the pair:
      \[(x - (2+3i))(x - (2-3i)) = x^2 - 4x + 13\]
    \item If $p$ also has real zero $x = 1$, the minimum degree is \textbf{3}.
  \end{itemize}
\end{frame}

% EVEN AND ODD FUNCTIONS
\begin{frame}[t, plain]
  \navyheader{Even and Odd Functions}
  \vspace{0.15cm}
  \begin{block}{Definitions and Algebraic Test}
    \begin{itemize}
      \item \textbf{Even}: $f(-x) = f(x)$ \quad $\Rightarrow$ symmetric about the $y$-axis
      \item \textbf{Odd}: $f(-x) = -f(x)$ \quad $\Rightarrow$ symmetric about the origin
      \item \textbf{Test}: substitute $-x$, simplify, compare to $f(x)$ and $-f(x)$
    \end{itemize}
  \end{block}
  \vspace{0.2cm}
  \renewcommand{\arraystretch}{1.5}
  \begin{tabular}{|l|l|l|}
  \hline
  \textbf{Function} & \textbf{$f(-x)$} & \textbf{Result} \\
  \hline
  $f(x) = x^4 - 2x^2$ & $x^4 - 2x^2 = f(x)$ & \textbf{Even} \\
  \hline
  $g(x) = x^3 - x$ & $-x^3 + x = -g(x)$ & \textbf{Odd} \\
  \hline
  $h(x) = x^2 + x$ & $x^2 - x$ (neither) & \textbf{Neither} \\
  \hline
  \end{tabular}
  \vspace{0.15cm}

  \textit{Caution: even/odd function $\neq$ even/odd degree!}
\end{frame}

% CLASS PRACTICE
\begin{frame}[t, plain]
  \navyheader{Class Practice}
  \vspace{0.3cm}
  Answer each question. Be ready to share your reasoning.
  \vspace{0.2cm}
  \begin{enumerate}
    \setlength{\itemsep}{10pt}
    \item A polynomial with real coefficients has zeros $x = 4$ (multiplicity 2) and
          $x = 1 - i$. List all zeros. What is the minimum degree?
    \item Determine whether $k(x) = 2x^6 - 3x^4 + x^2$ is even, odd, or neither.
          Show your algebraic test.
    \item For $p(x) = (x - 5)^3(x + 2)^2$, which zeros cause the graph to cross the
          $x$-axis, and which cause it to bounce?
  \end{enumerate}
\end{frame}

\end{document}
